% Source: https://artofproblemsolving.com/community/c3h1917446_question_on_lcms_and_running_laps
% Stable identity: aops-topic-1917446-q01
(a) Brooks gains on Avery
\[
\dfrac{1}{78}-\dfrac{1}{91}=\dfrac{91-78}{78\cdot91}=\dfrac{13}{7098}=\dfrac{1}{546}
\]
lap per second. He overtakes Avery again when he has gained one whole lap, after $\boxed{546}$ seconds.

(b) With lap times $a<b$ the overtaking time is $\dfrac{1}{\dfrac{1}{a}-\dfrac{1}{b}}=\dfrac{ab}{b-a}$, while $\operatorname{lcm}(a,b)=\dfrac{ab}{\gcd(a,b)}$. They agree exactly when $b-a=\gcd(a,b)$. Here $\gcd(78,91)=13=91-78$, so $\operatorname{lcm}(78,91)=546$ as well, by coincidence. For lap times $60$ and $100$ seconds the overtaking time is $\dfrac{6000}{40}=150$ seconds, but $\operatorname{lcm}(60,100)=300$: at $300$ seconds both runners are back at the start, but the faster one has already overtaken the slower one at $150$ seconds.
